\section{Постановка задачи}
	\tikzsetfigurename{FLPQ_ProblemState_}

Пусть нам дан конечный ориентированный помеченный граф $\mscrG = \langle V, E, L \rangle$.
Функция $\omega(\pi) = \omega((v_0, l_0, v_1), (v_1, l_1, v_2), \dots, (v_{n-1}, l_{n-1}, v_n)) = l_0 \cdot l_1 \cdot \dots \cdot l_{n-1}$ строит слово по пути посредством конкатенации меток рёбер вдоль этого пути.
Очевидно, для пустого пути данная функция будет возвращать пустое слово, а для пути длины $n  > 0$~--- непустое слово длины $n$.

Если теперь рассматривать задачу поиска путей, то окажется, что то множество путей, которое мы хотим найти, задаёт множество слов, то есть язык.
А значит, критерий поиска мы можем сформулировать следующим образом: нас интересуют такие пути, что слова, составленные из меток вдоль них, принадлежат заданному языку.

\begin{definition}[Задача поиска путей с ограничениями в терминах формальных языков]
    \label{def:FLPQ}
    \emph{Задача поиска путей с ограничениями в терминах формальных языков} заключается в поиске множества путей $\Pi = \{\pi \mid \omega(\pi) \in \mscrL\}$.
\end{definition}

В задаче поиска путей мы можем накладывать дополнительные ограничения на путь (например, чтобы он был простым, кратчайшим или Эйлеровым~\sidecite{kupferman2016eulerian}), но это уже другая история.

Другим вариантом постановки задачи является задача достижимости.

\begin{definition}[Задача достижимости]
    \label{def:Reachability}
    \emph{Задача достижимости} заключается в поиске множества пар вершин, для которых найдется путь с началом и концом в этих вершинах, что слово, составленное из меток рёбер пути, будет принадлежать заданному языку
    \[\Pi' = \{(v_{i}, v_{j}) \mid \exists v_{i} \pi v_{j}, \omega(\pi) \in \mscrL\}.\]
\end{definition}

При этом, множество $\Pi$ может являться бесконечным, тогда как $\Pi'$ конечно, по причине конечности графа $\mscrG$.

Язык $\mscrL$ может принадлежать разным классам и быть задан разными способами.
Например, он может быть регулярным, контекстно-свободным, или многокомпонентным контекстно-свободным.

Если $\mscrL$~--- регулярный, $\mscrG$ можно рассматривать как недетерминированный конечный автомат (НКА), в котором все вершины являются одновременно и стартовыми, и конечными.
Тогда задача поиска путей, в которой $\mscrL$~--- регулярный, сводится к пересечению двух регулярных языков.

Более подробно мы рассмотрим случай, когда $\mscrL$~--- контекстно-свободный язык.

Путь $G = \langle \Sigma, N, P \rangle$~--- контекстно-свободная грамматика.
Будем считать, что $L \subseteq \Sigma$.
Мы не фиксируем стартовый нетерминал в определении грамматики, поэтому, чтобы описать язык, задаваемый ей, нам необходимо отдельно зафиксировать стартовый нетерминал.
Таким образом, будем говорить, что $L(G,N_i) = \{ w \mid N_i \xRightarrow[G]{*} w  \}$~--- это язык задаваемый грамматикой $G$ со стартовым нетерминалом $N_i$.

\begin{example}
    Пример задачи поиска путей.

    Дана грамматика $G$, задающая язык $\mscrL = a^n b^n$:
    \begin{align*}
        S & \to a b   \\
        S & \to a S b
    \end{align*}
    И дан граф $\mscrG$, представленный на рисунке~\ref{fig:all_paths_example}.
\begin{marginfigure}
    \begin{center}
        \input{figures/edge_labelled_graph_a_3_loop_b_2_loop}
    \end{center}
    \caption{Входной граф}
    \label{fig:all_paths_example}
\end{marginfigure}

    Кратчайшими путями, принадлежащими множеству $\Pi = \{\pi \mid \omega(\pi) \in \mscrL\}$, являются, например, пути, представленные на изображении~\ref{fig:paths}.

\begin{marginfigure}
    \begin{center}
        \input{part_03_GraphAnalysis/chapter_10_FLPQ/figures/path1}
        \input{part_03_GraphAnalysis/chapter_10_FLPQ/figures/path2}
    \end{center}
    \caption{Примеры путей}
    \label{fig:paths}
\end{marginfigure}

\end{example}
